% ---------------------------------------------------------------------------
% Results fragment for Section 9 (bilateral partnership indices), subsection "Results"
% (subsec:bilateral_results). Numbers: paper/results/bilateral.json; figures and tables:
% scripts/paper/fig_bilateral.R (final rerun of 2026-10-12). The script rebuilds the Friendshore and
% Opportunity indices, the best-five summaries and the core-view PSI from the stored components with
% the pipeline's functions and stops unless they reproduce the stored tables.
% Figures: bilateral_fs_heatmap, bilateral_fs_components, bilateral_cmfdi_exposure,
% bilateral_psi_rank, bilateral_fs_po. Tables: tables/bilateral_coverage.tex,
% tables/bilateral_psi_views.tex.
% All correlations are Spearman rank correlations. Rankings of partners use the pipeline's
% rankability rule (sovereign, imputed weight share at most 0.5): psi_*_tbl rank / rankable.
% Citations needed but not in references.bib: the CM-FDI database (Net Zero Industrial Policy Lab,
% Johns Hopkins University) and the NEIS Allied Industrial Base working paper (October 2026); both
% are named in the prose without a citation command.
% ---------------------------------------------------------------------------

% Generated by scripts/paper/fig_bilateral.R; numbers in paper/results/bilateral.json (coverage).
\begin{table}[t]
\centering
\begin{threeparttable}
\caption{Coverage of the bilateral partnership tables}
\label{tab:bilateral_coverage}
\footnotesize
\begin{tabular}{@{}l r r r@{}}
\toprule
 & Friendshore & Opportunity & Development\\
\midrule
Reporters & 219 & 219 & --\\
Partners & 219 & 219 & 132\\
Reporter--partner pairs & 47{,}742 & 47{,}742 & --\\
Cells $(i,j,t,s)$ or $(j,t,s)$ & 1{,}534{,}002 & 1{,}342{,}942 & 4{,}356\\
Technology--stage cells with an HS6 basket & 18 & 18 & --\\
Reporters with observed bilateral trade & 208 & -- & --\\
Share of basket cells with observed trade & 77.0\% & 77.5\% & --\\
\quad of which from the partner's mirror declaration & 21.9\% & 21.8\% & --\\
Reporters with IMF outward positions (other cells) & 62 & -- & --\\
Share of cells with an investment component & 18.1\% & -- & --\\
\quad CM-FDI cells: share with a project tie & 7.2\% & -- & --\\
\quad CM-FDI cells: share with Chinese-parent exposure & 12.4\% & -- & --\\
Share of cells with reporter need missing & 12.5\% & -- & --\\
Partners with stage aid (Up/Mid/Down) & -- & -- & 121/131/131\\
Partners in $\PSI_j$ (core view); pairs in $\PSI_{ij}$ & \multicolumn{3}{c}{219 partners; 47{,}742 pairs}\\
\bottomrule
\end{tabular}
\begin{tablenotes}[flushleft]\footnotesize
\item \textit{Notes:} Final rerun of 2026-10-12. Trade shares are over the dyad cells of the 18 technology--stage cells with an HS6 basket; observed trade is the importer's own declaration or, where the importer does not report, the exporter's mirror declaration times the CIF/FOB factor 1.06 (all-reporter Comtrade dyads, 2020 and 2024). The investment component is the IMF outward direct-investment position of the reporter in the partner (16 conduit financial centres excluded) except in the four cells (Batteries, Electric Vehicles, Solar and Wind Midstream) where it is the reporter's cross-border manufacturing projects in the partner recorded in CM-FDI. Development Potential covers the partners that are not high-income. Source: \texttt{partner\_*\_tbl.rds}, \texttt{partner\_trade\_coverage\_tbl.rds}, \texttt{psi\_tbl.rds}, \texttt{psi\_dyad\_tbl.rds}; \texttt{scripts/paper/fig\_bilateral.R}.
\end{tablenotes}
\end{threeparttable}
\end{table}


\paragraph{Coverage and the trade matrix.} Table~\ref{tab:bilateral_coverage} reports the coverage
of the three pillars. Friendshore and Partnership Opportunity rate every ordered pair of the 219
economies, 47,742 pairs, in 33 and 30 technology--stage cells respectively; Development Potential
rates the 132 partners that are not high-income. The trade components rest on the all-reporter
Comtrade dyads for 2020 and 2024 \citep{uncomtrade}: 168 economies report imports in 2020 and 143 in
2024, and 208 economies appear as reporter or partner. A flow is taken from the importer's own
declaration where the importer reports, and otherwise from the exporter's mirror declaration times
the CIF/FOB factor 1.06. Over the dyad cells of the 18 technology--stage cells that have an HS6
basket (nine Upstream and nine Midstream; no Downstream stage has one), 77.0 percent of the
Friendshore trade components are observed, 17.4 percent as positive flows and 59.7 percent as
recorded zeros, and 23.0 percent are imputed at the cell mean. Mirror declarations supply 21.9
percent of the observed values; without them 39.9 percent of the trade components would be imputed.
The mirror rule matters for reporters that do not declare their own imports, which enter the matrix
only through their partners' declarations. It matters little for the core members, whose imports
are almost all self-declared (0.3 percent of their observed import components come from mirrors and
5.2 percent are imputed). Dropping the mirror fill and re-running every later step leaves the
headline ranking almost unchanged (Spearman 0.997 with the headline $\PSI_j$ over the 192 ranked
partners; median rank change 1, the same top fifteen), but it moves $C^{\PO}_j$ (Spearman 0.816),
because a third of the core exporters' observed export flows to their markets are mirrors of the
markets' own import declarations; the largest losers would be the United Arab Emirates (rank 91 to
126) and Russia (105 to 117).

The outward-investment component has two sources. In 29 of the 33 cells it is the reporter's IMF
outward direct-investment position in the partner, with 16 conduit financial centres (among them
the Netherlands, Luxembourg, Ireland, Singapore, Switzerland, Hong Kong and Cyprus) excluded as
counterparts and as reporters; 62 economies report positions, 23 of them core members, covering
9,171 pairs. In Batteries, Electric Vehicles, Solar and Wind Midstream it is instead the
reporter's cross-border manufacturing projects in the partner recorded in the CM-FDI database of
the Net Zero Industrial Policy Lab at Johns Hopkins University: 905 of its 1,042 projects, those
in manufacturing that are announced, under construction, completed or operating (cancelled, paused
and closed plants are dropped), give 234 investor--destination
pairs from 29 investing economies, 18 of them core members. Because the ties are sparse, 7.2 percent
of the dyad cells in these four sectors carry an investment component, against 19.5 percent of the
cells elsewhere. Reporter need is missing in 12.5 percent of the Friendshore cells, among them every cell of
Electric Grid Downstream and the two Geothermal stages. Of the 219 partners, 192 are ranked in the
headline view: 24 are not sovereign, and Liechtenstein, Monaco and Kosovo have more than half of
their index weight imputed. The median imputed weight share of a Friendshore dyad is 0.42 over all
reporters and 0.26 over the core members, whose trade and investment data are the most complete.

\paragraph{The coalition Friendshore matrix.} Figure~\ref{fig:bilateral_fs_heatmap} shows
$\FS_{ijts}$ among the 31 core members in Batteries Midstream and Solar Midstream, with China as a
further supplier. Every one of the 930 core dyads in each cell has an observed import component.
Supplier identity matters more than importer identity: the standard deviation of the suppliers'
mean ratings is 0.052 and 0.051, that of the importers' 0.036 and 0.033. The matrix is strongly
asymmetric. The Spearman correlation of $\FS_{ijts}$ with $\FS_{jits}$ over the 465 unordered core
pairs is $-0.40$ in Batteries Midstream and $-0.41$ in Solar Midstream, and $-0.14$ on average
within cells. The asymmetry is built into the index: the reporter's need and the partner's
capability enter with positive weight, so a capable, secure economy rates its partners low on
their capability and is itself rated high, and the reverse holds for a needy importer. China is
rated by the core members at 0.613 in Batteries Midstream and 0.605 in Solar Midstream on average,
last among the 32 suppliers shown in both cells, although its trade component (0.96 and 0.94) and
capability (1.00) are the highest. The alignment component, zero between China and every core member, and the
Chinese-parent component, zero for China by construction, account for the gap: without the alignment component China's mean rating would be 0.825 and 0.818.

\begin{figure}[tbp]
  \centering
  \includegraphics[width=\textwidth]{bilateral_fs_heatmap}
  \caption{Friendshore index among the core members, (a) Batteries Midstream and (b) Solar
  Midstream. Rows are importers $i$, columns suppliers $j$; colour is $\FS_{ijts}$
  (equation~\eqref{eq:fs_index}), the within-cell S-curve of the weighted mean of the seven
  components less the climate penalty. Rows and columns are the 31 core coalition members ordered
  by their mean rating as suppliers; the last column is China as a supplier to each core member.
  The diagonal is empty. Every dyad shown has an observed import component (importer's declaration
  or mirror). Source: \texttt{partner\_friendshore\_tbl.rds}.}
  \label{fig:bilateral_fs_heatmap}
\end{figure}

\paragraph{What makes a top supplier.} Figure~\ref{fig:bilateral_fs_components} decomposes the
Friendshore composite of the ten highest-ranked suppliers of the United States and of Japan, over
the five cells behind each dyad's best-five summary $T^{\FS}_{ij}$. Suppliers are ranked among the
partners that are sovereign and whose Friendshore dyad has at most half of its weight imputed (158
partners for the United States, 170 for Japan). The leading suppliers of the United States are the
United Kingdom, Germany, Canada, the Netherlands, Ireland, Australia, France, Slovakia, Estonia and
Lithuania, and those of Japan the United States, Portugal, Austria, Poland, Canada, Ireland,
Thailand, Czechia, Estonia and Germany; all are U.S.-leaning except Thailand. Alignment is the
largest single contribution to their composites (29.1 percent for the United States' top ten and
29.0 percent for Japan's), followed by partner capability (23.0 and 24.9 percent), imports (18.7 and
20.1 percent) and reporter need (18.7 and 16.9 percent); outward investment contributes 9.9 and 6.9
percent and the Chinese-parent component 0.7 and 2.1 percent, because few of these dyads' best cells
are among the four cells the component covers. The climate penalty is small (0.018 and 0.033 on
average). Half of the best cells of these dyads have no HS6 basket, so the ranking of top suppliers
rests largely on the components that are defined everywhere: need, capability and alignment. China
ranks 70th among the suppliers of the United States and 75th among those of Japan ($T^{\FS}=0.71$);
without the alignment component it would rank 10th and 8th. None of the ten largest outward-investment
counterparts of the two reporters is a conduit centre, since these are excluded; the United States'
largest are the United Kingdom, Australia, Taiwan, Malaysia and Canada.

\begin{figure}[tbp]
  \centering
  \includegraphics[width=\textwidth]{bilateral_fs_components}
  \caption{Composition of the Friendshore composite for the ten highest-ranked suppliers of the
  United States (left) and Japan (right). Each bar is the mean, over the five technology--stage
  cells behind the dyad's best-five summary $T^{\FS}_{ij}$ (in parentheses), of the weighted
  component contributions $w_c x_c/\sum_{c'\,\mathrm{obs}}w_{c'}$, which sum to the pre-penalty
  composite; the climate penalty is the negative segment. Components and weights: imports from $j$
  (2), reporter need (2), partner capability (2), outward investment in $j$ (2; IMF positions, or
  CM-FDI projects in Batteries, Electric Vehicles, Solar and Wind Midstream), UNGA bloc alignment
  (2) and Chinese-parent exposure scored as $1-$share (1, defined only in those four cells); the
  zero-weight reporter capability term is omitted. Suppliers are ranked among partners that are
  sovereign and whose Friendshore dyad has an imputed weight share of at most 0.5. Source:
  \texttt{partner\_friendshore\_tbl.rds}, \texttt{partner\_friendshore\_top5\_tbl.rds}; penalty
  recomputed from the technology emission index and the Climate Action Tracker ratings.}
  \label{fig:bilateral_fs_components}
\end{figure}

\paragraph{The alignment component.} UNGA bloc alignment carries weight 2 of a total of 11, 18.2
percent when every component is present, but a larger share of the composite in practice: on
average 0.125, or 24.1 percent of the pre-penalty composite. For a core reporter its contribution
is 0.231 for a U.S.-leaning partner, 0.115 for a nonaligned partner and zero for a China-leaning
one. Setting its weight to zero and re-running every later step changes the headline ranking more
than any other single choice examined here: the Spearman correlation of $\PSI_j$ with the headline
is 0.869 and that of $C^{\FS}_j$ 0.716, the median partner moves 17 places and the largest move is
81. China would rise from 44th to 3rd, India from 39th to 16th, Indonesia from 62nd to 15th, Viet
Nam from 64th to 20th, Russia from 105th to 47th, the United Arab Emirates from 91st to 21st and
Oman from 140th to 59th; 19 rather than 24 core members would be in the top 25. Among the top
fifteen, China, Saudi Arabia and Indonesia would replace Sweden, Korea and Ireland. In the headline
the core members occupy 24 of the top 25 places; the first non-member is Slovenia (25th), a member of
the EU node. The headline ranking is therefore in large part a statement about political alignment
with the United States, by construction, and readers who prefer a purely economic reading of
supplier strength should use the variant without it.

\paragraph{Chinese-parent exposure.} The seventh Friendshore component measures how much of a
supplier's foreign-invested manufacturing in a technology is owned by parents headquartered in
China, Hong Kong or Macao, from the 905 CM-FDI projects used, 406 of which have such a parent. The
share blends capital-weighted shares among projects with disclosed capital and project counts
among the rest, attributes half of a joint venture to each parent, and is defined only where the
supplier hosts at least two foreign-invested projects in the technology; China, Hong Kong and Macao
take the value one. It is defined for 24 to 28 suppliers per covered cell, 47 economies in all, 31
of them coalition members (Figure~\ref{fig:bilateral_cmfdi_exposure}). Over the 94 supplier--technology
cells outside China, Hong Kong and Macao the median share is 0.50; 48 cells are at 0.5 or above and
19 at zero. Exposure is highest in solar (median 0.75) and lowest in wind (median 0). Several
partners whose supply the coalition relies on are heavily exposed: Türkiye (batteries 1.00, solar
0.97, electric vehicles 0.82), Thailand (batteries 0.97, electric vehicles 0.88), Malaysia
(electric vehicles 0.90), Viet Nam (electric vehicles 1.00), Indonesia (solar 0.85) and Saudi
Arabia and Oman (solar 1.00), and among the core Korea (batteries 0.91, wind 1.00), Hungary
(electric vehicles 0.64, batteries 0.55) and Slovakia (electric vehicles 0.61). The United States
(at most 0.27, in solar), Canada and Mexico (zero in batteries) are among the least exposed.

The component operationalises one of the upgrades that the NEIS working paper on the Allied
Industrial Base proposes for indices of this kind: two suppliers in two countries do not diversify
a buyer's supply if both depend on the same parent company. Its weight is deliberately modest and
its coverage narrow, so its effect on the headline is small: the component contributes 0.051 to
the composite on average where it is defined (8.1 percent), the Spearman correlation of $\FS$ with
and without it within the four cells is 0.99, and the headline $\PSI_j$ ranking is unchanged up to
four places (Hungary would rise from 20th to 16th without it; Denmark and Finland would fall three
places). Its effect is concentrated where it should be. In the core members' ratings of a supplier
in a covered cell, the change in $\FS$ due to the component correlates at $-0.93$ with the
Chinese-parent share; suppliers with a share of at least one half lose 0.063 on average, those
below 0.2 gain 0.043, and China loses between 0.085 and 0.095 in each of the four cells. The
measure captures ownership of foreign-invested plants only; it does not see domestic firms'
dependence on Chinese equipment, precursors or licences, the other channels of common-mode risk.

\begin{figure}[tbp]
  \centering
  \includegraphics[width=\textwidth]{bilateral_cmfdi_exposure}
  \caption{Chinese-parent ownership exposure of suppliers. Panel (a): for each economy hosting at
  least two foreign-invested manufacturing projects in a technology, the share of those projects
  owned by parents headquartered in China, Hong Kong or Macao (capital-weighted where capital is
  disclosed, by count otherwise; joint ventures attribute half to each parent), in the Midstream
  cells of batteries, electric vehicles, solar and wind; grey cells are undefined (fewer than two
  projects). Rows are grouped by coalition tier and ordered by the mean share; China, Hong Kong and
  Macao, set to one by construction, are omitted. Panel (b): the mean Friendshore rating of the
  economy by the 31 core members over its covered cells, with the component (filled) and with its
  weight set to zero and every later step re-run (open). Source: CM-FDI database, Net Zero
  Industrial Policy Lab, Johns Hopkins University (CC BY 4.0), 905 of 1,042 projects;
  \texttt{partner\_cm\_fdi\_exposure\_tbl.rds}, \texttt{partner\_friendshore\_tbl.rds}.}
  \label{fig:bilateral_cmfdi_exposure}
\end{figure}

\paragraph{The partner-level PSI.} Figure~\ref{fig:bilateral_psi_rank} reports the headline
$\PSI_j$, averaged over the 31 core members' views, and its three category scores. The United
States (0.936), Germany (0.934), the Netherlands, the United Kingdom, France, Italy, Spain,
Belgium, Poland and Czechia (0.904) are the ten highest; the core members hold 24 of the top 25
places (median rank 16; New Zealand is 41st and Ukraine 60th), and the first non-members are
mostly EU members outside the core, which enter the network through the EU node: Slovenia (25th),
Lithuania, Malta, Estonia, Latvia and Cyprus, with North Macedonia among them. Partners rank lower (median 64.5), from Saudi Arabia (38th), India (39th) and Türkiye (43rd)
to Ghana, Oman and Gabon (139th to 141st). Among non-members China ranks 44th, Taiwan 47th and
Switzerland 51st. China is a strong market and a weak supplier in this view: it ranks 21st on
$C^{\PO}_j$ and 57th on $C^{\FS}_j$. Taiwan is 8th on $C^{\PO}_j$ and 55th on $C^{\FS}_j$. The
dispersion of the category scores differs widely: the standard deviation across ranked partners
is 0.190 for $C^{\FS}_j$, 0.031 for $C^{\PO}_j$ and 0.190 for $C^{\DP}_j$, so the Opportunity
score, nominally 40 percent of the index, accounts for only 5.5 percent of the variance of
$\PSI_j$ among the 130 partners with all three categories (Friendshore 63.8 percent, Development
30.7 percent) and 14.5 percent among the 62 without Development. The categories carry distinct
information: their Spearman correlations are 0.59 ($C^{\FS}$ and $C^{\PO}$, 192 partners), 0.31
($C^{\FS}$ and $C^{\DP}$) and 0.29 ($C^{\PO}$ and $C^{\DP}$, 130 partners). High-income partners
average 0.827 and the other partners on the World Bank list 0.672. The market score follows the core members' observed exports to the market (Spearman 0.82) and
correlates negatively with the market's own security need (Spearman $-0.54$), although need enters
it with a positive weight: the economies most in need of supply are mostly small importers. The strongest supplier--weak market contrasts are small
European states and some low-income suppliers (San Marino, Andorra, Albania, Chad, Niger); the
reverse are markets such as the United Arab Emirates (17th on $C^{\PO}_j$, 132nd on $C^{\FS}_j$),
Kazakhstan, Liberia and Uganda.

\begin{figure}[tbp]
  \centering
  \includegraphics[width=\textwidth]{bilateral_psi_rank}
  \caption{Partnership Strength Index in the headline view. Diamonds are $\PSI_j$
  (equation~\eqref{eq:psi}) with $C^{\FS}_j$ and $C^{\PO}_j$ averaged over the 31 core coalition
  members as reporters; dots are the three category scores, Friendshore $C^{\FS}_j$, Partnership
  Opportunity $C^{\PO}_j$ and Development Potential $C^{\DP}_j$ (partners that are not high-income
  only). Panels: the 31 core members, the 26 partners, the ten highest-ranked non-members with
  China, Taiwan and Switzerland, and the five highest-scoring partners that are not ranked. Ranks
  in parentheses are among the 192 rankable partners (sovereign, imputed weight share at most 0.5);
  open diamonds and ``n.r.'' mark partners without a rank. Source: \texttt{psi\_tbl.rds}.}
  \label{fig:bilateral_psi_rank}
\end{figure}

\paragraph{Views of the index.} The headline averages the categories over the core members. Two
alternatives average over the full coalition (57 reporters) and over every reporter (the world
view, 218 per partner); 192, 192 and 188 partners are ranked in the three views
(Table~\ref{tab:bilateral_psi_views}). The core and full-coalition views agree closely (Spearman
0.981 over 192 partners, median rank change 5, 13 of the top fifteen in common); the core and world
views less so (0.881 over 188, median change 17, ten in common), and the full-coalition and world
views lie in between (0.948). The views differ mainly in how they place the large non-aligned and
China-leaning economies. China is 44th in the core view, 9th in the full-coalition view and first in
the world view; India 39th, 15th and 2nd; Indonesia 62nd, 32nd and 6th; Viet Nam 64th, 46th and
12th. The world view's top fifteen add China, India, Indonesia, Viet Nam and Malaysia and drop
Belgium, Poland, Czechia, Sweden and Ireland. The partners whose rating depends on who is asking
are thus exactly those whose alignment differs from that of the core; the choice of the core
members' view as the headline is a choice to rate partners from the vantage point of the economies
that anchor the coalition.

% Generated by scripts/paper/fig_bilateral.R; numbers in paper/results/bilateral.json (psi_views).
\begin{table}[t]
\centering
\begin{threeparttable}
\caption{The fifteen highest-ranked partners in the three views of the Partnership Strength Index}
\label{tab:bilateral_psi_views}
\footnotesize
\begin{tabular}{@{}r l S[table-format=1.3] l S[table-format=1.3] l S[table-format=1.3]@{}}
\toprule
 & \multicolumn{2}{c}{Core members (31 reporters)} & \multicolumn{2}{c}{Full coalition (57)} & \multicolumn{2}{c}{World (218)}\\
\cmidrule(lr){2-3}\cmidrule(lr){4-5}\cmidrule(l){6-7}
Rank & Partner & {$\PSI_j$} & Partner & {$\PSI_j$} & Partner & {$\PSI_j$}\\
\midrule
1 & United States & 0.936 & United States & 0.893 & China & 0.813\\
2 & Germany & 0.934 & Germany & 0.879 & India & 0.787\\
3 & Netherlands & 0.922 & France & 0.868 & France & 0.783\\
4 & United Kingdom & 0.919 & Netherlands & 0.864 & United States & 0.782\\
5 & France & 0.917 & Spain & 0.861 & Germany & 0.780\\
6 & Italy & 0.914 & United Kingdom & 0.857 & Indonesia & 0.768\\
7 & Spain & 0.908 & Italy & 0.854 & Netherlands & 0.764\\
8 & Belgium & 0.906 & South Korea & 0.852 & South Korea & 0.761\\
9 & Poland & 0.906 & China & 0.850 & Spain & 0.760\\
10 & Czechia & 0.904 & Belgium & 0.848 & United Kingdom & 0.757\\
11 & Austria & 0.900 & Canada & 0.847 & Italy & 0.755\\
12 & Sweden & 0.896 & Poland & 0.840 & Viet Nam & 0.753\\
13 & Canada & 0.896 & Czechia & 0.838 & Canada & 0.750\\
14 & South Korea & 0.894 & Austria & 0.838 & Austria & 0.746\\
15 & Ireland & 0.893 & India & 0.836 & Malaysia & 0.746\\
\bottomrule
\end{tabular}
\begin{tablenotes}[flushleft]\footnotesize
\item \textit{Notes:} $\PSI_j$ (equation~\eqref{eq:psi}) with $C^{\FS}_j$ and $C^{\PO}_j$ averaged over the reporters of each view: the 31 core coalition members (the headline), the 57 core and partner-tier members, and every other reporter (218 per partner). Ranks among the rankable partners of each view (sovereign, imputed weight share at most 0.5): 192, 192 and 188 partners. Spearman correlations of $\PSI_j$ across the partners ranked in both views: core and full coalition 0.981 (192), core and world 0.881 (188), full coalition and world 0.948 (188). Source: \texttt{psi\_tbl.rds}, \texttt{psi\_full\_coalition\_tbl.rds}, \texttt{psi\_world\_tbl.rds}; \texttt{scripts/paper/fig\_bilateral.R}.
\end{tablenotes}
\end{threeparttable}
\end{table}


\paragraph{Supplier and market views of the same edge.} Figure~\ref{fig:bilateral_fs_po} compares,
for the coalition, the importer's Friendshore rating of a supplier, $\FS_{cp}$, with the
supplier's Opportunity rating of that importer as a market, $\PO_{pc}$. The two pillars rate the
same supply edge from its two ends and share inputs: the partner capability term of Friendshore and
the reporter capability term of Opportunity are the same Economic Opportunity cell (within-cell
Spearman 1.0), and the imports of $c$ from $p$ and the exports of $p$ to $c$ are the same flow
where both are observed (0.81). They agree to a degree, with a cell-level Spearman of 0.715 over
95,312 matched coalition cells (0.764 among core members), 0.733 for the mean over the matched cells
of the 3,192 ordered coalition pairs and 0.703 for the best-five summaries; given the shared
capability term, the mean within-cell partial correlation is 0.44. Alignment separates them, since it
enters Friendshore only: for pairs in the same bloc Friendshore exceeds Opportunity by 0.157 on
average, for pairs in opposite blocs it falls short by 0.24, and for the 2,070 pairs with a
nonaligned member the two agree on average (0.016). The pairs the market view favours most are
pairs in opposite blocs, such as Luxembourg supplying Oman ($\FS=0.37$ against $\PO=0.71$) and
Indonesia supplying Canada (0.42 against 0.75); those the supplier view favours most have Ukraine as
the supplier, to Greece (0.70 against 0.23) and Sweden. At the dyad level the coalition pairs with
the highest $\PSI_{ij}$ are Canada's rating of the United States (0.981), Spain's of Portugal, the
United States' of the United Kingdom, and Austria's and Czechia's of Germany (0.970--0.976).

\begin{figure}[tbp]
  \centering
  \includegraphics[width=\textwidth]{bilateral_fs_po}
  \caption{Supplier and market ratings of the same supply edge, coalition members. Horizontal:
  $\FS_{cp}$, member $c$'s Friendshore rating of member $p$ as a supplier; vertical: $\PO_{pc}$,
  $p$'s Opportunity rating of $c$ as a market. Panel (a): Batteries Midstream, one point per ordered
  pair of the 57 members (3,192 pairs). Panel (b): mean over the technology--stage cells rated by
  both pillars (28 to 30 per pair). Colour is the pair's dyadic UNGA alignment (1 same bloc, 0.5 when
  either is nonaligned, 0 opposite blocs); the line is the diagonal; the three pairs most favoured by
  each view are labelled. Source: \texttt{partner\_friendshore\_tbl.rds},
  \texttt{partner\_opportunity\_tbl.rds}.}
  \label{fig:bilateral_fs_po}
\end{figure}

\paragraph{Development potential.} Among the 125 partners ranked on Development Potential, the
highest best-five summaries belong to Georgia, Mongolia and Mauritius (1.00), Tuvalu, China
(0.919), Serbia, Rwanda, Viet Nam, North Macedonia and Mauritania (0.872). Their best cells are
spread over the stages (40 percent Upstream, 26 percent Midstream, 34 percent Downstream), and the
two aid components, 35 percent of the nominal weight, account for 43.7 percent of the composite in
those cells. China's high score is a reminder that the pillar rates a partner's development
indicators and capacity to absorb supply-chain investment, not its alignment.
